\chapter{Conclusions and Future Work}
\label{ch:conclusions}

\section{Conclusions}

This project asked a simple question: does adding a semantic signal to an acoustic one, and
then deriving a cost-sensitive decision from the result, improve end-of-turn detection for a
French medical voice agent? On a clean public benchmark the answer is yes. Fusion raises the
AUC from 0.814, for the acoustic model alone, to 0.897, a statistically significant gain of 0.083, and neither base model is fine-tuned in that
pipeline. A medically fine-tuned lexical model exists and serves the in-domain path and the
live demonstration, but it is not what produces this number.

What the project learned is worth as much as what it built. Fusion only helps when the
lexical model is in distribution: applied to generic French, a medically tuned model
degrades the decision below the baseline it was meant to improve. Calibration goes
unreported, though a calibrated probability is the precondition for deciding from a stated
cost rather than an arbitrary timer: every external detector evaluated here is miscalibrated
on this data and none publishes any measurement of it. And the off-the-shelf pipeline does
not transfer to production, for reasons Section~\ref{sec:gold} sets out and O4 returns to.

There is one more finding I did not expect, and it cuts against the framing of the project.
Section~\ref{sec:hardcase} shows the fusion gaining least on precisely the cases the work
exists for, and shows that generic French barely contains those cases at all. The motivating
problem is real, but the public benchmark cannot measure it, and the lexical signal is
weakest exactly where I argued it would be strongest.

This dissertation contributes a calibration-first, cost-aware method for composing existing
models, tested in a real setting, and it measures where that method stops working.

\section{Evaluation against the objectives}

\textbf{O1}, a modular real-time pipeline whose decision never blocks on the transcript, is
met: the package is in place, covered by 45 automated tests, and demonstrated end to end
from the command line and in the browser.

\textbf{O2}, fusing off-the-shelf acoustic and lexical models and testing whether this beats
the acoustic model alone, is met (Section~\ref{sec:fusionresults}).

\textbf{O3}, calibrating the probability and deciding from a cost, is met as analysis and
not as a property of the running system. Calibration was measured across the whole
leaderboard, and the operating point is derived from the cost ratio rather than from a fixed
delay, but that derivation lives in the evaluation: the shipped controller still maps
confidence onto a hand-set silence ladder.

\textbf{O4}, instrumenting production and validating in domain, was carried out, and its
result is negative. The capture and the gold pipeline work, with 709 calls captured and 636 scored, and a validated
converter, but the validation showed that the off-the-shelf pipeline does not transfer: the
acoustic branch is at chance on 8\,kHz audio and the lexical branch plateaus near 0.60 on
weak labels. The objective was delivered as engineering and as measurement; the hypothesis
it was built to test was not confirmed.

Several limitations bound what has been shown. On generic French the method plateaus around
0.897, and passing that needs larger models, which is outside the contribution. The system
trails the commercial streaming services on latency, because it waits for a micro-pause
while they commit as the text arrives. The in-domain labels are weak, and the French
benchmark subset is small at roughly a thousand silences.

There are things I would do differently. I would instrument production in the first week rather
than the seventh, so that a properly annotated in-domain set existed in time instead of a
weakly supervised one. I would check the audio bandwidth of the target deployment before choosing an acoustic
model. The 16\,kHz-to-8\,kHz gap surfaced at validation, which is far too late for it to
surface.
And I would run the per-case breakdown early: had I known in week three that the public
benchmark contains six spelling cases, I would have known from the start that it could not
test the thing the project was for.

The three pillars of this work were set out before implementation began, in the critical
review described in Chapter~\ref{ch:intro}, and two of them survived intact. The third
bent: I evaluate at the micro-pause rather than continuously as speech arrives, which is
exactly the latency gap of Section~\ref{sec:fusionresults} and the first item of future
work below. Being overtaken mid-project was the other thing that shaped the contribution. Narrowing the
claim to calibration and cost, instead of competing on fusion itself, was, I still think,
the right response.

\section{Future work}

The first direction is latency, which is what stands between this system and deployment. It
decides at the micro-pause, which is why it trails services that commit as the text arrives.
Two levers exist: a streaming lexical model that scores word by word and can decide before
the pause, and the stability of the recogniser's interim hypotheses, how quickly the
transcript stops changing, as an early cue. That second signal is only observable in the
partial-hypothesis trajectories captured in production, and cannot be tested on the public
benchmark, which supplies final words only.

The second answers the in-domain result directly: adapt the acoustic model to the 8\,kHz
telephone band, by retraining or by starting from a natively narrowband model, to remove the
saturation, and hand-annotate a few hundred pauses to obtain clean end-of-turn labels and a
trustworthy in-domain measurement, which weak supervision cannot provide.

The third is a sequential optimal-stopping formulation. The current rule applies a myopic
threshold to a calibrated probability. At high cost ratios the swept oracle turns
conservative through timing rather than through the threshold, and a myopic rule cannot
match that. A multi-step cost model, in the tradition of \citet{wald1945sequential}, would
track the oracle across every cost regime.
